\section{Latihan}
\label{sec:latihan-bab9}

\begin{enumerate}
    \item \textbf{Konseptual.} Jelaskan perbedaan instruksi mesin dan
          direktif assembler, serta mengapa direktif disebut
          \eng{pseudo-op}. Berikan dua contoh direktif yang mengisi
          memori dan dua contoh yang tidak mempengaruhi isi berkas
          \file{.exe}.

    \item \textbf{Konseptual.} Sebutkan direktif data yang tepat
          untuk setiap kebutuhan berikut, lalu tuliskan
          deklarasinya: (a) satu nilai \hex{0A} untuk pindah baris;
          (b) sebuah tabel berisi delapan nilai 16-bit untuk
          \hex{0000}; (c) sebuah penampung 32 byte yang diisi saat
          program berjalan; (d) sebuah string \code{'SEMARANG'}
          yang siap dicetak dengan layanan DOS fungsi \hex{09}.

    \item \textbf{Analisis.} Diberikan blok data berikut. Hitung
          total byte yang disediakan dan tuliskan \eng{offset}
          setiap label.

\begin{lstlisting}[language=MASM, numbers=none]
.DATA
    a   DB 'TEST'
    b   DW 2 DUP(0)
    c   DD 1, 2
    d   DB 3 DUP(?)
    e   DT 0
\end{lstlisting}

    \item \textbf{Analisis \eng{little-endian}.} Tuliskan urutan byte
          di memori untuk deklarasi \lstasm|DW 0102h| dan
          \lstasm|DD 03040506h|. Jelaskan mengapa urutannya
          demikian, lalu tentukan nilai yang terbaca pada
          \reg{AX} bila program menjalankan \lstasm|MOV AX, w|
          dengan \lstasm|w DW 0102h|.

    \item \textbf{Analisis kesalahan.} Perhatikan potongan berikut.

\begin{lstlisting}[language=MASM, numbers=none]
BARIS   EQU 25
KOLOM   EQU 80
BARIS   EQU 50
\end{lstlisting}

          \begin{enumerate}
              \item Baris mana yang menimbulkan kesalahan, dan apa
                    pesan yang dikeluarkan TASM?
              \item Ubah potongan itu agar sah, tetapi tetap
                    menghasilkan \code{BARIS} bernilai \dec{50}
                    pada pemakaian terakhir.
              \item Jelaskan mengapa pemakaian \code{=} tidak
                    menimbulkan masalah pada kasus ini.
          \end{enumerate}

    \item \textbf{Analisis alur translasi.} Sebuah berkas sumber
          berakhir dengan \lstasm|END| tanpa nama label.
          \begin{enumerate}
              \item Apakah TASM melaporkan kesalahan? Jelaskan.
              \item Apakah TLINK melaporkan kesalahan? Jelaskan
                    pesan yang muncul.
              \item Bagaimana cara memperbaikinya?
          \end{enumerate}

    \item \textbf{Penulisan program.} Tulislah program lengkap yang
          mendeklarasikan sebuah tabel berisi lima nilai
          \hex{41}--\hex{45} dengan \dir{DB}, mencetaknya dengan
          layanan DOS fungsi \hex{02} satu karakter per
          pemanggilan, lalu keluar dengan fungsi \hex{4C}.

    \item \textbf{Penulisan program dan pembacaan berkas hasil.}
          Deklarasikan semua nomor fungsi DOS yang dipakai dengan
          \dir{EQU}, lalu rakit dengan \file{tasm /l} dan
          \file{tlink /m}. Salin tabel segmen dari berkas
          \file{.map}, lalu jelaskan hubungan besar segmen
          \file{\_DATA} dengan jumlah byte deklarasi data pada
          program Anda.
\end{enumerate}

\begin{refleksibox}
Jawab pertanyaan berikut pada jurnal praktikum Anda.
\begin{enumerate}
    \item Direktif mana yang paling sering Anda lupakan fungsinya?
          Bagaimana Anda mengingatnya kembali?
    \item Setelah membaca berkas \file{.lst} dan \file{.map},
          apakah gambaran Anda tentang kerja assembler dan linker
          menjadi lebih jelas? Bagian mana yang paling menolong?
    \item Kesalahan perakitan apa yang Anda temui pada latihan ini,
          dan apa pelajaran yang Anda ambil dari kesalahan
          tersebut?
\end{enumerate}
\end{refleksibox}
